%=== spec/fonctions/lectgraph_11.tex || Fonctions || Graphique

\hfill\begin{GraphiqueTikz}[x=0.75cm,y=0.75cm,Xmin=-3,Xmax=3,Ymin=-2.25,Ymax=1.75]
  \TracerAxesGrilles[Origine,Grille=false]{}{}
  \GenererPolynomeLagrange{-2.5,-1.45,0.25,2.55}{-1.9,0.85,0.75,-0.25}
  \TracerCourbe[Couleur=red]{polylagrange(x)}
  \PlacerTexte[Couleur=red]{(0.95,0.4)}{$\mathcal{C}$}
  \MarquerPts[Couleur=blue]{(-2.5,-1.9)/$A$/left,(-1.45,0.85)/$B$/above left,(0.25,0.75)/$R$/above,(2.55,-0.25)/$S$/below right}
\end{GraphiqueTikz}\hfill\null

\smallskip

On a représenté ci-dessus la courbe $\mathcal{C}$ d'une fonction $f$. Les points $A$, $B$, $R$ et $S$ appartiennent à la courbe $\mathcal{C}$. Leurs abscisses sont notées respectivement $x_{A}$, $x_{B}$, $x_{R}$ et $x_{S}$.

\smallskip

L'inéquation $x \times f(x) > 0$ est vérifiée par :

\smallskip

\eamreponsesautom[2]{$x_{A}$ et $x_{B}$}{$x_{A}$ et $x_{R}$}{$x_{A}$ et $x_{S}$}{$x_{A}$, $x_{B}$ et $x_{S}$}